\chapter{Conclusion and Future Work}
\label{ch:conclusion}

\section{Answers to the Research Questions}

\paragraph{RQ1 -- Does flow-based ML generalise across time?}
The answer depends on the task formulation. Multi-class classification
does not generalise: macro F1 collapses from 0.86 on the random
stratified split to 0.44 on the day-based temporal split, with all four
models within 0.04 of one another. The collapse is structural --- Friday's
classes are not in training --- and model tuning cannot recover the loss.
Binary detection, however, does largely generalise: random forest reaches
F1 = 0.86 on the day split and ROC-AUC = $0.926 \pm 0.053$ across
leave-one-day-out folds.

\paragraph{RQ2 -- Are flow-based detectors robust to low-effort evasion?}
No. The undefended random forest flips 21.4\% of attack flows at
$\varepsilon = 0.25\sigma$ on the perturbable feature subset, with a
looser headline evasion rate of 24.4\% when the attack is run unbounded,
and 83.6\% of successful evasions move a single feature (Flow IAT Min).
Adversarial flows transfer to XGBoost at 86.4\%, so ensemble-based
defences are ineffective against this kind of attack. Naïve adversarial
training fails: clean F1 drops from 0.86 to 0.73 while evasion increases
to 26.8\%.

\paragraph{RQ3 -- Can model output drive triage?}
Yes, via static MITRE ATT\&CK mapping. The pipeline produces structured
per-flow alerts annotated with technique ID, kill-chain phase, and a
one-sentence justification, ready for direct routing to SOC playbooks.
Sub-class attribution on the day split is unreliable for the same
structural reason that drives the multi-class collapse, so the deployed
system defaults to ``unspecified attack'' when the multi-class prediction
is out-of-distribution.

\section{Future Work}

Five directions stand out. First, cross-dataset LODO --- training on
CICIDS days and testing on UNSW splits --- would extend the temporal
generalisation claim across capture environments. Second, the static
MITRE mapping could be replaced with a learned mapping using richer
telemetry. Third, a stronger adversarial training defence (PGD with
curriculum scheduling, balanced augmentation, classifier-rebalanced
thresholds) would test whether the negative result on naïve adversarial
training can be overturned with a more careful approach. Fourth,
streaming evaluation --- deploying as a Kafka consumer with continuous
concept-drift detection --- would test the deployment artefacts under
realistic operational load. Fifth, modern tabular and sequence models
(TabNet, FT-Transformer, LSTM-over-IATs) deserve a careful comparison
against the tree-ensemble baseline.

\section{Closing Statement}

Flow-based machine learning for intrusion detection works in practice,
but only when it is evaluated under protocols that respect the temporal
structure of the data. On the two benchmarks studied here, headline
numbers from random stratified splits substantially overstate the
deployment-relevant performance reported under day-based and LODO
splits. Binary detection is the more credible deployment framing on
these benchmarks, MITRE-mapped triage is operationally valuable, and
the entire pipeline is reproducible in 45 to 60 minutes on a laptop.
The integration --- honest temporal evaluation, adversarial robustness,
and operational MITRE triage in one reproducible pipeline --- is what
distinguishes this thesis from a benchmark replay. The unsolved problem
is adversarial robustness: any deployed detector should assume that a
determined adversary will eventually evade it, and the right response is
to combine the machine-learning detector with network-layer mitigations,
host-based detection, and human analysts in a defence-in-depth
architecture.
